% ---------------------------------------------------------------------
\section{单举 $B$ 衰变}\label{sec:heavy-inclusive}

单举衰变率与 $R_{e^+e^-}$、DIS 一样是编时乘积前向矩阵元的吸收部分，可按 $1/m_b$ 作 OPE；在谱的端点，强子末态是 $p_X^2\sim m_b\Lqcd$ 的喷注，领头的一串项须求和成形状函数（shape function）。对 $\bar B\to X_s\gamma$，$p_b=m_bv+k$，光子 $q=\sqrt2E_\gamma n_B$ 沿 $-z$，$p_X=m_bv+k-q$ 沿 $+z$，
\begin{equation}\label{eq:heavy-pX2}
p_X^2=m_b\,\big(m_b-2E_\gamma+\omega\big)+\order{\Lqcd^2},\qquad \omega\equiv\frac{n_A\cdot k}{n_A\cdot v}=k^0-k^3 ,
\end{equation}
末态只通过 $\omega$ 感受 $b$ 夸克的剩余动量（脚本 \texttt{s3\_\allowbreak inclusive.py}）。

\begin{theorem}[单举 $B$ 衰变的重夸克展开 \tagA]\label{thm:heavy-inclusive-b}
(a)（局域 OPE）总半轻与总辐射衰变率，以及在强子末态能量上以宽度 $\Delta\gg\Lqcd$ 的光滑权重平均的谱，有 $1/m_b$ 的渐近展开：领头项是自由 $b$ 夸克衰变（含其 $\as$ 修正），没有 $\order{\Lqcd/m_b}$ 项，$1/m_b^2$ 项由
\begin{equation}\label{eq:heavy-lambda12}
\lambda_1=\frac{\langle\bar B|\bar h_v(\ii D)^2h_v|\bar B\rangle}{2m_B},\qquad
\lambda_2=\frac{\langle\bar B|\bar h_v\frac g2\sigma_{\mu\nu}G^{\mu\nu}h_v|\bar B\rangle}{6m_B}
\end{equation}
决定（$[\ii D^\mu,\ii D^\nu]=\ii gG^{\mu\nu}$；$m_{B^*}^2-m_B^2=4\lambda_2+\order{\Lqcd^3/m_b}$）。例如，以极点质量 $m_b$ 表示，
\begin{equation}\label{eq:heavy-Gamma-bu}
\Gamma(\bar B\to X_u\ell\bar\nu)=\frac{G_F^2|V_{ub}|^2m_b^5}{192\pi^3}\Big[1+\order{\as}+\frac{\lambda_1-9\lambda_2}{2m_b^2}+\order{\frac{\as\Lqcd^2}{m_b^2},\frac{\Lqcd^3}{m_b^3}}\Big].
\end{equation}
(b)（端点区）当 $m_b-2E_\gamma=\order{\Lqcd}$ 时，在 $\Lqcd/m_b$ 的领头幂次、对 $\as$ 的所有阶，
\begin{equation}\label{eq:heavy-HJS}
\frac{\dd\Gamma(\bar B\to X_s\gamma)}{\dd E_\gamma}=2\Gamma_0\,H(m_b,\mu)\int_{2E_\gamma-m_b}^{\bar\Lambda}\dd\omega\;m_b\,J\big(m_b(m_b-2E_\gamma+\omega),\mu\big)\,S(\omega,\mu),
\end{equation}
其中 $\Gamma_0$ 是 $b\to s\gamma$ 的领头阶宽度，$H=1+\order{\as}$，$J(p^2,\mu)=\delta(p^2)+\order{\as}$ 是沿 $n_A$ 的夸克喷注函数，$\bar\Lambda=m_B-m_b$，形状函数
\begin{equation}\label{eq:heavy-shape}
S(\omega)=\frac{1}{2m_B}\Big\langle\bar B(v)\Big|\,\bar h_v\,\delta\Big(\omega-\frac{n_A\cdot\ii D}{n_A\cdot v}\Big)h_v\Big|\bar B(v)\Big\rangle
\end{equation}
（双局域形式中两个 $h_v$ 之间是沿 $n_A$ 的类光 Wilson 线）；$\bar B\to X_u\ell\bar\nu$ 的端点区有同一形状函数的类似因子化。
非微扰地位：\cref{asm:intro-wightman}、\cref{asm:intro-composite}、\cref{asm:intro-af-coeff}（MA0、MA4）；端点（形状函数）区的因子化另需\cref{asm:intro-bridge}（MA1、MA2）。
% errtype: E-S
\uses{thm:dis-ope, thm:scet-decoupling, asm:intro-wightman, asm:intro-composite, asm:intro-af-coeff, asm:intro-bridge}
\end{theorem}

\begin{proofsketch}
(1) 对编时乘积作 OPE，其全阶地位同\cref{thm:dis-ope}（[A]）。(2) 领头算符 $\bar b\gamma^0b$ 由 $b$ 数守恒精确归一；量纲 4 的唯一候选 $\bar h_v\,\ii v\cdot D\,h_v$ 由 HQET 运动方程为零，故没有 $1/m_b$ 修正\cite{ChayGeorgiGrinstein1990,BigiUraltsevVainshtein1992}；量纲 5 的动能与色磁算符给出 $\lambda_{1,2}$。把 Manohar--Wise 的 $\Gamma/\Gamma_b=1+5E_b-6K_b-2G_b$\cite[式~(5.7)，$\rho\to0$]{ManoharWise1994}中的 $K_b=-\frac{\lambda_1}{2m_b^2}$、$G_b=-\frac{3\lambda_2}{2m_b^2}$（他们取 $[D_\mu,D_\nu]=\ii gG_{\mu\nu}$）代入即得\cref{eq:heavy-Gamma-bu}。(3) OPE 只在离开物理割线处可靠，围道在小强子质量处靠近割线，故 (a) 只对总宽度与涂抹的谱成立（与\cref{ch:ee}类似）。(4) 端点区中 $[\Lqcd/(m_b-2E_\gamma)]^n$ 型的项同阶，求和成形状函数\cite{Neubert1994a,Neubert1994b,BigiEtAl1994}。在 SCET 中，$b\to s\gamma$ 流匹配到 SCET 流（$H$ 为匹配系数的模方）；$B$ 中动量 $\sim\Lqcd$ 的组分相对于喷注是超软的，经 BPS 场重定义退耦（\cref{thm:scet-decoupling}），留下沿 $n_A$ 的 Wilson 线，即\cref{eq:heavy-shape}；喷注部分给出 $J$。RG 一致性使 Sudakov 对数可以求和\cite{BauerFlemingLuke2000}；图形层次上的同一结构见\cite{KorchemskySterman1994}，喷注函数见\cite{Bauer:2003pi,Becher:2006qw}（E-S 型陈述）。
\end{proofsketch}

\paragraph{地位与缺口}
全阶论证（[A]）：(a) 的地位是微扰 OPE 的地位（\cref{thm:dis-ope}，[A]），(b) 的地位是 SCET 因子化的地位（\cref{rem:scet-assumptions}，模式完备性，E-S 型陈述）。用于\emph{固定运动学点}上的衰变率需要局域对偶（\cref{conj:ee-duality}，[C]），这不是上述定理的前提，而是其应用所需的额外输入。极点质量的红外重整化子带来相对量级 $\Lqcd/m_b$ 的歧义，与系数级数的发散相抵消；改用短距离质量后歧义消失，与没有 $1/m_b$ 项相符\cite{Bigi:1994em,Beneke:1994sw}。升级途径：(a) 随微扰 OPE 升级；(b) 需证明端点区模式的完备性。

\paragraph{树图检验。}
取 $H=1$、$J=\delta(p^2)$ 得 $\frac{\dd\Gamma}{\dd E_\gamma}=2\Gamma_0S(2E_\gamma-m_b)$；由 $\int S=1$、$\int\omega S=0$、$\int\omega^2S=-\lambda_1/3$（树图层次），$\langle E_\gamma\rangle=m_b/2$，谱的方差为 $-\lambda_1/12$，支撑 $\omega\le\bar\Lambda$\cite{Neubert1994b}把端点移到 $m_B/2$；$\lambda_1$ 项的系数 $\frac12$ 是时间膨胀 $\langle m_b/E_b\rangle$（脚本 \texttt{s7\_\allowbreak extra.py}）。
